\section{Постановка задачи исследования}
\label{sec:problem}

Инфраструктурная боль, из которой выросла работа, двоякая. С одной
стороны, логика построения таблицы \textit{IP-префикс~$\rightarrow$~ASN}
жила внутри BIFIT COLLECTOR, тогда как BIFIT MITIGATOR нуждался в тех же
данных: без отдельного сервиса неизбежны дублирование и расхождение
копий. С другой стороны, backend \texttt{ranger} обновлялся по таймеру
полной заменой снимка; одиночная вставка сводилась к перестройке всей
структуры, так что поток BGP-событий нельзя было применять к in-memory
базе по одному событию без катастрофической цены.

Сервис, рассматриваемый в работе, обслуживает несколько внутренних
потребителей и должен отвечать на три типа запросов:

\begin{itemize}
\item к какой автономной системе относится конкретный IP-адрес;
\item каков полный текущий снимок таблицы соответствия префиксов
автономным системам;
\item какие изменения в таблице произошли с момента подписки клиента.
\end{itemize}

Прежняя связка «встроенная логика + \texttt{ranger}» удовлетворяла
части требований к чтению, но не позволяла ни разделить продукты по
границе сервиса, ни сделать потоковое обновление базовым режимом при лаге порядка
секунд.

\subsection{Исследовательская гипотеза}

Опираясь на свойства BGP (глава~2) и критерии выбора in-memory структуры
(глава~3), сформулирую гипотезу так:

\begin{quote}
\textit{если вынести функцию сопоставления
\textit{IP-префикс~$\rightarrow$~ASN} из состава одного продукта в
автономный микросервис с общим gRPC API для нескольких потребителей,
перейти от периодической полной замены таблицы к потоковому применению
BGP-событий к in-memory базе по схеме «начальный снимок + поток
обновлений» и реализовать внутреннее хранилище на основе шардированного
Adaptive~Radix~Tree, то таблицу можно поддерживать актуальной в реальном
времени и снизить стоимость одного инкрементального обновления на
несколько порядков по сравнению с пакетной моделью на \texttt{ranger}.}
\end{quote}

Архитектурный шаблон «снимок + поток» -- частный случай более общей
stateful-streaming парадигмы, описанной у Акидау и
соавторов~\cite{akidau2015dataflow}; он хорошо сочетается с
лог-ориентированными платформами Kafka~\cite{kreps2011kafka} и
Flink~\cite{carbone2015flink} и согласуется с трактовкой состояния как
функции последовательности событий поверх базового
снимка~\cite{kleppmann2017ddia}.

\subsection{Функциональные требования}

\begin{enumerate}
\item Сервис должен уметь подключаться к источнику маршрутного
состояния (узлу GoBGP) и считывать полный снимок IPv4 и IPv6
маршрутов.
\item Сервис должен уметь принимать поток update-событий и отражать
его во внутреннем состоянии.
\item Сервис должен предоставлять API для точечного lookup по
IP-адресу.
\item Сервис должен предоставлять API для получения полного снимка
таблицы.
\item Сервис должен предоставлять API подписки на инкрементальные
изменения.
\item Сервис должен корректно обрабатывать вставки, удаления и
повторное построение состояния после ресинхронизации.
\end{enumerate}

\subsection{Нефункциональные требования}

\begin{enumerate}
\item Одиночное обновление не должно требовать полной перестройки всей
таблицы.
\item Чтение из таблицы должно оставаться достаточно быстрым для
прикладного использования.
\item Полная выгрузка может быть медленнее, чем в batch-оптимизированном
решении, но должна оставаться практически приемлемой.
\item Архитектура должна позволять одновременно обслуживать нескольких
клиентов через единый API.
\item Реализация должна быть достаточно прозрачной и сопровождаемой,
включая возможность профилирования и наблюдения внутреннего состояния.
\end{enumerate}

\subsection{Критерии успеха}

Критерии, по которым результат работы должен признаваться успешным,
вытекают из исследовательской гипотезы и требований:

\begin{enumerate}
\item стоимость отражения одного изменения маршрута уменьшена не
менее чем на несколько порядков;
\item время точечного lookup сохранено в прикладно-приемлемых
пределах;
\item возможность построения полного дампа сохранена и остаётся
работоспособной для целевого размера таблиц;
\item потоковая модель работает как основной операционный сценарий, а
не как факультативное дополнение к batch-режиму.
\end{enumerate}

%==============================================================
